\section{Теоремы}

\begin{theorem}[О целочисленности потока]
	Если пропускные способности всех дуг конечной сети целочисленны, то существует максимальный целочисленный поток, который может быть найден методом расстановки пометок за конечное число итераций.
\end{theorem}

\begin{theorem}[Форда-Фалкерсона]
	В любой конечной сети величина максимального потока равна пропускной способности минимального разреза.
\end{theorem}

\begin{theorem}[Критерий разрешимости ЗТЗ]
	ЗТС разрешима тогда и только тогда, когда выполнено условие баланса:
	\[
		\sum_{i=1}^{m} a_i = \sum_{j=1}^{n} b_j.
	\]
\end{theorem}

\begin{theorem}[Критерий базисности плана ЗТЗ]
	План ЗТЗ является базисным, если и только если из его основных коммуникаций нельзя составить цикл.
\end{theorem}

\begin{theorem}[Критерий оптимальности плана ЗТЗ]
	План $X = (x_{ij} )$ ЗТЗ является оптимальным $\iff \exists $ потенциалы $u_i,v_j: \ \forall i=\overline{1,m} , \ j = \overline{1,n} $:
	\begin{enumerate}
		\item $x_{ij} > 0 \implies v_j - u_i = c_{ij} $.
		\item $x_{ij} = 0 \implies v_j - u_i \leqslant c_{ij} $.
	\end{enumerate}
\end{theorem}

\begin{theorem}[О целочисленности плана ЗТЗ]
	Если все объемы производства $a_i$ и потребления $b_j$ -- целые числа, то методом потенциалов будет найден целочисленный оптимальный план ЗТЗ.
\end{theorem}

\begin{theorem}[Свойства ситуаций равновесия]
	Пусть $(x^1,y^1),(x^2,y^2)$ -- две различные ситуации равновесия в антагонистической игре. Тогда:
	\begin{enumerate}
		\item $H(x^1,y^1) = H(x^2,y^2)$.
		\item $(x^1,y^2),(x^2,y^1)$ тоже являются ситуациями равновесия.
	\end{enumerate}
\end{theorem}

\begin{theorem}[Критерий существования седловой точки]
  Платежная функция имеет седловую точку, если и только если $u^0 = v^0$.
\end{theorem}

\begin{theorem}[О минимаксе]
  Для любой матричной игры имеет место равенство $u^* = v^*$.
\end{theorem}

\begin{theorem}[УБ "назад"]
	\[
		\forall x \in X, \ L(x) = \underset{y\in A(x)}{\max} \big\{f(x,y) + L(y)\big\},
	\]
	\[
		L(x_k) = a_k, \ \forall x_k \in X_{\beta } .
	\]
\end{theorem}

\begin{theorem}[УБ "вперед"]
	\[
		\forall x \in X, \ R(x) = \underset{y\in B(x)}{\max} \big\{f(y,x) + R(y)\big\},
	\]
	\[
		R(x_k) = b_k, \ \forall x_k \in X_{\alpha } .
	\]
\end{theorem}

\begin{theorem}
	В любой непротиворечивой сетевой модели существует хотя бы один критический путь.
\end{theorem}
